\section*{Kapitel 7 --- Komponenttest genom stubben}
\addcontentsline{toc}{section}{Kapitel 7 --- Komponenttest genom stubben}

\solution{ex:7:1}{Skriv buggen, se testet fånga den}
Uppgiften ber om en kodändring; frågorna besvaras här.
\ans{a} Fallen som kontrollerar den ekade framens innehåll och räknaren fallerar --- framen som
ekas har fel \code{dlc}, och fallet som anropar \code{run()} två gånger på en injicerad frame
avslöjar att tillståndet konsumerats. Fallet utan injicerad frame är fortfarande grönt, eftersom
\code{receive()} returnerar \code{false} redan första gången.
\ans{b} Nej, inget av dem. \code{Stub} och \code{Frame} gör exakt vad de lovar; ingen av dem är
trasig.
\ans{c} Ett enhetstest kan inte se en bugg som bor \emph{mellan} enheterna. Det krävs ett test som
kör flera delar tillsammans för att upptäcka att den ena använder den andra fel.

\solution{ex:7:2}{Lagergränsen}
\ans{a} Nej. \code{EchoNode} måste känna till interfacet --- det är dess enda beroende.
\ans{b} Nej. Testet håller i den konkreta typen; det är där \code{inject()} finns.
\ans{c} Ja. Implementationsfilen ligger också ovanför interfacet, och får inte veta vilken
drivrutin den kör mot.
\ans{d} Nej, den bryter inte mot gränsen --- \code{hasError()} och \code{clearError()} tillhör
interfacet. Men det är ändå tveksamt: \code{EchoNode} kan inte göra något meningsfullt åt ett fel
den inte kan tolka, och felflaggan är en bit för fyra orsaker. Felhanteringen hör hemma hos den
som äger drivrutinen och känner systemets krav, alltså i factoryn eller i lagret ovanför.
\ans{e} Ja, grovt. Casten namnger en konkret drivrutin \emph{och} når en privat metod i den. Den
skulle dessutom krascha så fort en \code{Stub} injicerades.

\solution{ex:7:3}{Ägande}
\ans{a} Det kompilerar inte. \code{Interface} är abstrakt och går inte att instansiera som värde;
kompilatorn säger att typen är abstrakt. Vore den det inte skulle medlemmen \emph{skiva} objektet
--- kopiera basklassdelen och kasta resten --- vilket är värre än ett kompileringsfel.
\ans{b} Det kompilerar. Referensen pekar då på ett förstört objekt, och \code{run()} är odefinierat
beteende. Kompilatorn kan inte se det; det är den ena kostnaden för att komponenten inte äger sin
drivrutin.
\ans{c} Därför att \code{EchoNode} inte äger den. Drivrutinens livstid är längre än komponentens,
den kan delas med andra komponenter, och ägandet ligger hos factoryn. En \code{unique_ptr} hade
påstått motsatsen, och en \code{shared_ptr} hade lagt till räkning som ingen behöver.

\solution{ex:7:4}{Läs sviten}
\ans{a} Fallet som bara kontrollerar den ekade framens \emph{innehåll} --- rätt \code{id},
\code{dlc} och data --- skulle vara grönt, eftersom ekot fortfarande sker. Det är därför sviten
också kontrollerar räknaren.
\ans{b} Trafik matas in med \code{stub.inject()} och läses tillbaka med \code{stub.receive()}.
Ingen av dem finns i \code{Interface}: \code{inject()} kan \code{Spi} inte implementera, och
\code{receive()} används här i en annan roll än interfacets --- testet inspekterar vad komponenten
\emph{skickade}, vilket stubben råkar lagra i samma fält.
\ans{c} Att \code{run()} konsumerar framen: efter första anropet finns ingenting kvar att eka, så
räknaren ska stå på 1 och inte 2. Det är kravet som fångar dubbelanropet i övning~\ref{ex:7:1}.
